%  This is a example presentation using the provided beamertheme
%  <jacknjo> compiled with pdflatex.
%
%  Copyright (C) 2016  Florian Hofmann - JacknJo
%
%  This file is part of the beamertemplate jacknjo which 
%  can be found at <https://github.com/JacknJo/jacksbeamertheme>.
%
%  The theme and this example is free software: you can
%  redistribute it and/or modify it under the terms of the GNU 
%  General Public License as published by the Free Software 
%  Foundation, either version 2 of the License, or (at your 
%  option) any later version.
%
%  The theme is distributed in the hope that it will be useful,
%  but WITHOUT ANY WARRANTY; without even the implied warranty of
%  MERCHANTABILITY or FITNESS FOR A PARTICULAR PURPOSE.  See the
%  GNU General Public License for more details.
%
%  You should have received a copy of the GNU General Public 
%  License along with this theme. 
%  If not, see <http://www.gnu.org/licenses/>.

\documentclass[
	12pt, 				% fontsize; available are 10pt, 11pt, 12pt
	t,					% define to start from top of the page
	aspectratio=169,	% page aspectratio; available are 43, 169, 54, 32, 1610, 149
	%handout			% define if you only want a handout version without the overlay makros
	]{beamer}
	
% ##### choose theme and define it's color #####
\usetheme{jacknjo}
\colortheme{1} % [0] blue; [1] orange; [2] green; [3] red; [4] turquoise;

% ##### language and encoding #####
\usepackage[T1]{fontenc}
\usepackage[ngerman]{babel} % by using the babel package, date shows as dd Month yyyy

% ##### other packages you might need #####
\usepackage{lipsum}
%\usepackage{minted}		% Can be used for source code highlighting. Mention the workaround in file <var/mintedExample.tex>

\usepackage{datetime}
\ddmmyyyydate \renewcommand{\dateseparator}{.}


% The useLualatex command only matters if you want to change the font. If you don't want to just leave it 0. 
% Setting useLualatex to 1 forces you to compile your project with the lualatex compiler. 
% Also one of the following \setmainfont{fontname} lines has to be uncommented 
% (notice that the .ttf or .otf font has to be installed on your system)

\newcommand{\useLualatex}{0} % [0] pdflatex; [1] lualatex
\ifnum\useLualatex=0
	\usepackage[utf8]{inputenc}
\else
	\usepackage{fontspec}
	%X##### simple font, straight letters, mathmode - ok
	%\setmainfont[Mapping=tex-text]{Lucida Sans} 
	
	%X##### similar to ubuntu, less bold mathmode - ok
	%\setmainfont[Mapping=tex-text]{Trebuchet MS}
	
	%X##### bold text, not for math mode
	%\setmainfont[Mapping=tex-text]{Ubuntu}  % Ubuntu Light as alternative
	
	%X#####
	%\setmainfont[Mapping=tex-text]{Liberation Sans}
	
	%X##### great looking, favorite
	%\setmainfont[Mapping=tex-text]{Open Sans} 
	
	%X##### good looking, strict and standard font
	%\setmainfont[Mapping=tex-text]{Fira Sans} 
	
	%X##### great looking, playful round & easy to read, no bold font!
	%\setmainfont[Mapping=tex-text]{ABeeZee}\renewcommand{\textbf}[1]{\textit{#1}}
	
	%X##### playful font, not so serious
	%\setmainfont[Mapping=tex-text]{Josefin Sans SemiBold} 
	
	\let\sfdefault\rmdefault
\fi


% ##### custom environments #####
\newenvironment{myjustify}{\begin{justify}\vspace*{-0.475cm}}{\end{justify}}

% ##### custom commands #####
\newcommand{\pictureat}[4]{%
	\only<+(1)->{
		\begin{textblock*}{\paperwidth}(#1, #2)%
			\includegraphics[#3]{#4}%
		\end{textblock*}}}
		
% This custom command can be used with the theme and provides a zoom in of the specified picture
% it has to be specified manually by the first command if width = \pagewidth or height = \pageheight
% to fit on the screen. The second command specifies the path to the picture.
\newcommand{\fillpic}[2]{
	\clearpage
	\begin{textblock*}{\paperwidth}(0pt,0pt)
		\begin{tikzpicture}
		\useasboundingbox (0, 0) rectangle(\paperwidth,\paperheight);
		\fill [color = black, fill opacity=0.75](0, 0) rectangle(\paperwidth,\paperheight);
		\node (label) at (0.5\paperwidth,0.5\paperheight){\includegraphics[#1]{#2}};
		\end{tikzpicture}
	\end{textblock*}
}

% ##### Titlepage settings - positive and negative lengths are accepted #####
\setlength{\titleoffset}{0.5cm}
\setlength{\subtitleoffset}{0.35cm}
\setlength{\dateoffset}{0cm}

% ##### Frame settings - positive and negative lengths are accepted #####
\setlength{\framesubtitleoffset}{0pt}
\setlength{\frametitleoffset}{0.05cm}
\setlength{\frametextoffset}{0.1cm}
\setlength{\framefootlineoffset}{0cm}

% ##### Title page Information ##### 
\title{Jack's beamer theme}
\subtitle{My very first own \LaTeX{} design}
\author{%
		Florian Hofmann\\
		alias JacknJo
		}
\date{Friedrichshafen\\\today}

\newcommand{\middleText}{\copyright{} your copyright notice}
\newcommand{\outlineSection}{Outline}
\newcommand{\biblographySection}{Literature}

% ##### Use this line to place your logo. {posx}{posy}{options}{filename}
\renewcommand{\uselogo}{1} %uncommented -> logo on ; commented -> logo off
    \setlength{\logox}{\paperwidth  - 1.25cm}
    \setlength{\logoy}{\paperheight - 0.65cm}
    \newcommand{\logographics}{\includegraphics[width=2cm]{figures/jackssolutions}}
    
% ##### Add your literature.bib here #####
\addbibresource{var/literature.bib}


\begin{document}
		
	\begin{frame}[noframenumbering]
		\titlepage
	\end{frame}


	\begin{frame}{Outline}
		\tableofcontents
	\end{frame}


	\outlineFrame{Introduction}
	
	
	\begin{frame}{Background}{\& personal inspiration}
		Hello all of you! 
		
		I'm Jack, 20 years old and student in electrical engineering in Friedrichshafen in Germany.
		This is my first try to design a own \LaTeX{} template.
		
		It's advantages are:
		\begin{itemize}
			\item clean structure
			\item nice \& consistent design
			\item easy to use with a small basic \LaTeX{} knowledge
			\item useful additional commands
			\item export file is \texttt{.pdf} which can be opened nearly anywhere!
		\end{itemize}
	\end{frame}
	
	
	\outlineFrame{Utilities}
	
	
	\begin{frame}{Text}{}
		Sometimes you need to present a paragraph of text, which looks better justified.
		\vfill
		\begin{myjustify}
			\lipsum[2]
		\end{myjustify}
		\vfill
	\end{frame}
	
	
	\begin{frame}{Itemize}{}
		One of the most common ways to create slides is to use lots of itemizes.
		
		They are used for:
		\begin{itemize}
			\item better readability
			\item as short descriptions
			\begin{itemize}
				\item for easy points
				\item or difficult
				\begin{itemize}
					\item topics
					\item statements
				\end{itemize}
			\end{itemize}
			\item <2-> they also can be made visible item by item
			\item <3-> like this!
		\end{itemize}
	\end{frame}
	
	
	\begin{frame}{Todo}{list of all tasks}
		Many projects have deadlines and tasks you have to fulfill by a certain date.
		Presenting them the following way is very easy and leaves a good and professional impression.
		\vfill
		\begin{tabularx}{\textwidth}{Xll}
			\textbf{task} & \textbf{deadline} & \textbf{state}\\
			basic tutorials in creating a beamer template & 01.02.03 & \myCheck \\
			starting a own development & 01.03.16 & \myCheck \\
			creating a presentation with all the theme's capabilities & 03.04.16 & \myCheck \\
			finishing the beamer template & 01.06.16 & \myUnCheck\\
		\end{tabularx}
		\vfill
	\end{frame}
	
	
	\begin{frame}{Math environment}{presenting formula is very easy in \LaTeX}
		\vfill
		\begin{equation*}
			f(x)=\sum_{i=0}^\infty \frac{f^{(i)}(x_0)}{i!}(x-x_0)^i
		\end{equation*}
		\vfill
		\begin{equation*}
			\displaystyle\frac{\pi^2}{6}=\lim_{n \to \infty}\sum_{k=1}^n \frac{1}{k^2}
		\end{equation*}
		\vfill
	\end{frame}
	
	
	\begin{frame}{Blocks}{example, alert and standard}
	    \begin{exampleblock}{Exampleblock}
			Inside this box you can place an example or a task for your audience:
			
			What is the square root of 144? $\sqrt{144} = $
		\end{exampleblock}
		\begin{alertblock}{Alertblock}			
			There is no real solution to a negative square root!
		\end{alertblock}
		\begin{block}{Block}
			Use this for any other reason you can imagine!
		\end{block}
	\end{frame}
	
	
	\begin{frame}{Comparison}{of advantages and disadvantages}
		This is a example comparison:
		\begin{center}
		\begin{tabularx}{0.8\textwidth}{X|X}
			\textbf{Advantages}
			 \begin{itemize}[<2->]
			 	\item lifetime
			 	\item easy handling
			 	\item costs
			 \end{itemize} &
			\textbf{Disadvantages}
			 \begin{itemize}[<3->]
			 	\item weight
			 	\item resolution
			 	\item integration time
			 \end{itemize}
		\end{tabularx}
		\end{center}
	\end{frame}
	
	
	\begin{frame}{Comparison}{of two items}
		This is a example comparison:
		\begin{center}
		\begin{tabularx}{\textwidth}{XcX}
			\textbf{Advantages} & & \textbf{Disadvantages}\\
			lifetime & \drawopp & weight \\
			easy handling & \drawopp & resolution \\
			costs & \drawopp & integration time (also works for long texts)\\
			another & \drawopp & ... \\
		\end{tabularx}
		\end{center}
	\end{frame}	
	
	
	\begin{frame}{Cites}{that get automatically listed in the last frame}
		The following randomly picked text about the \textbf{Day of the Barricades} is from wikipedia and listed in \cite{art:dayofB}:
		\glqq In the French Wars of Religion, the Day of the Barricades (in French: Journée des barricades), 12 May 1588, was an outwardly spontaneous public uprising in staunchly Catholic Paris against the moderate [...] \grqq.
		
		\glqq Horse racing is an equestrian performance sport, involving two or more jockeys riding horses over a set distance for competition. It is one of the most ancient of all sports and its basic premise – to identify which of two or more horses is the fastest over a set course or distance – has remained unchanged since the earliest times. \grqq{} is a text about \textbf{horse racing} provided by \cite{art:horseRacing}
	\end{frame}
	
	
	\outlineFrame{Graphics}	
	
	
	\begin{frame}{Zoom}{stretch a picture to screen size}
		This is actually really usefull if you have a detailed picture and need the space to explain it's content.
		\vfill
		\begin{tabular}{ll}
			\includegraphics[width = 0.3\textwidth]{figures/img1} & \only<3->{
				\raisebox{0.15\textwidth}{
				\begin{minipage}{0.5\textwidth}
					\begin{itemize}
						\item Conclusion xyz
						\item and another more detailed one
						\begin{itemize}
							\item with this and
							\item this content
						\end{itemize}
						\item just edit the counter of the only command
					\end{itemize}
				\end{minipage}}
				} \\	
		\end{tabular}
		\only<2>{\fillpic{height=\paperheight}{figures/img1}}
		\vfill
	\end{frame}	
		
		
	\begin{frame}{Graphics}{absolute position}
		\begin{myjustify}
			\lipsum[3]
		\end{myjustify}
		\pictureat{1cm}{2.5cm}{width = 4cm}{figures/img1}
		\pictureat{7.25cm}{2.5cm}{width = 2cm}{figures/img1}
		\pictureat{7.25cm}{5.75cm}{width = 2cm}{figures/img1}
	\end{frame}
	
	
	\begin{frame}{Graphics}
	\framesubtitle{relative position in a tabular}
	This may be usefull for some iconified graphics.
		\begin{center}%
		\begin{tabular}{cc}%
			\visible<+(1)->{\includegraphics[width = 0.12\textwidth]{figures/img1}}&
			\visible<+(1)->{\includegraphics[width = 0.18\textwidth]{figures/img1}} \\
			\visible<+(1)->{\includegraphics[width = 0.14\textwidth]{figures/img1}} &
			\visible<+(1)->{\includegraphics[width = 0.14\textwidth]{figures/img1}} \\
		\end{tabular}%
		\end{center}%
	\end{frame}
	
	
	% To use minted in a beamer theme you have to use the following workaround and compile with the -shell-escape flag!
	%	\definecolor{bg}{rgb}{0.95,0.95,0.95}
	\newminted{python}{fontsize=\scriptsize, 
		linenos,
		numbersep=8pt,
		gobble=4,
		frame=lines,
		bgcolor=bg,
		framesep=3mm} 

	% check here: http://nochair.net/posts/2011/05-05-fragile-latex-beamer.html
	\defverbatim[colored]\exampleCode{
		\begin{pythoncode}
			import numpy as np
			import pylab as pl
			
			def f_x(x):
			return np.exp(x)+x**2-5*x
			
			def approx_f(x):
			return 1 -4*x +3./2*x**2
			
			xvals = np.arange(-4,4,0.1)
			fx_vals = [f_x(x) for x in xvals]
			
			pl.show()
		\end{pythoncode}
	}

	\begin{frame}{Code example}{with \textbackslash usepackage\{minted\}}
		\exampleCode
		\pause
		Overlays work!
	\end{frame}
	
	\section*{The End}
	\begin{frame}
		\vspace{2.2cm}
		\begin{center}
			\huge Thank you for watching this presentation.\\
			Have fun creating your own!
		\end{center}
		\vfill
	\end{frame}
	
	% ##### don't forget to compile with biblatex #####
	\printbibliographyframe
	
\end{document}